\section{Atividades}

\question[Atividade 1]{Varredura de vulnerabilidades no alvo}

\begin{figure}[H]
  \centering
  \includegraphics[width=0.99\textwidth]{./fig/fig1.png}
  \caption{Varredura da ferramenta ZAP 2.17.0 na página
  http://testasp.vulnweb.com}
  \label{fig:fig1}
\end{figure}

A atividade iniciou com acesso remoto via \textit{Cliente RDP}.

Foi utilizado o IP \texttt{192.168.98.12} com usuário \texttt{aluno} e senha
\texttt{rnpesr}.

Após o acesso, foi aberta a ferramenta \textit{OWASP ZAP} pelo menu:

\begin{itemize}
  \item \textbf{03 - Web Application Analysis}
  \item \textbf{Web Application Proxies}
  \item \textbf{zap}
\end{itemize}

Na primeira execução, os \textit{Add-ons} foram atualizados pela opção
\textit{Update All}.

Em seguida, foi iniciado um \textit{Automated Scan} com os parâmetros:

\begin{itemize}
  \item URL: \texttt{http://testasp.vulnweb.com/}
  \item \textit{Use traditional spider}: habilitado
  \item \textit{Use Ajax spider}: habilitado
  \item Agente: \textit{Firefox Headless}
\end{itemize}

O scan foi iniciado pelo botão \textit{Attack}.

O progresso foi acompanhado pela guia \textit{Active Scan} e os resultados
pelas seção \textit{Alertas}.

Ao final da análise, o OWASP ZAP realizou o mapeamento da aplicação e
identificou vulnerabilidades automaticamente.

As vulnerabilidades apresentaram descrição, impacto, severidade e recomendações
de mitigação.

\newpage

\question[Atividade 2]{Associação da metodologia OWASP Top 10 2021}

\begin{figure}[H]
  \centering
  \includegraphics[width=0.99\textwidth]{./fig/fig2.png}
  \caption{Vulnerabilidades associadas ao Top 10 do OWASP 2021 presentes no
  relatório gerado pela ferramente ZAP 2.17.0 na página
  http://testasp.vulnweb.com}
  \label{fig:fig2}
\end{figure}

Com base nos resultados do scan realizado no \textit{OWASP ZAP}, foram
identificadas quatro vulnerabilidades relacionadas ao \textit{OWASP Top 10
2021}.

\section*{Cross Site Scripting (Reflected)}

\textbf{Exploração resumida:}

A vulnerabilidade foi encontrada na rota \texttt{/showthread.asp?id=29}. Um
atacante pode enviar código JavaScript malicioso no parâmetro \texttt{tfText},
que é refletido pela aplicação sem tratamento adequado.

\textbf{Impacto provável:}

Roubo de sessão, execução de JavaScript arbitrário, captura de cookies e
redirecionamento do usuário para páginas maliciosas.

\textbf{Evidência mínima:}

\begin{itemize}
  \item Alerta ZAP: Cross Site Scripting (Reflected)
  \item CWE: CWE-79
  \item OWASP: A03:2021 - Injection
  \item URL:
    \texttt{http://testasp.vulnweb.com/showthread.asp?id=29}
\end{itemize}

\section*{SQL Injection}

\textbf{Exploração resumida:}

A aplicação aceita entradas sem validação adequada em consultas SQL. Um
atacante pode inserir comandos SQL maliciosos para manipular o banco de dados
da aplicação.

\textbf{Impacto provável:}

Leitura, alteração ou exclusão de dados sensíveis, além de possível
comprometimento completo do banco de dados.

\textbf{Evidência mínima:}

\begin{itemize}
  \item Alerta ZAP: SQL Injection
  \item OWASP: A03:2021 - Injection
  \item URL: \texttt{http://testasp.vulnweb.com/Search.asp?tfSearch=ZAP\%27}
\end{itemize}

\section*{Path Traversal}

\textbf{Exploração resumida:}

A vulnerabilidade permite manipular caminhos de arquivos enviados à aplicação.
Um atacante pode utilizar sequências como \texttt{../} para acessar diretórios
internos do servidor.

\textbf{Impacto provável:}

Leitura de arquivos sensíveis, exposição de configurações internas e possível
obtenção de credenciais do sistema.

\textbf{Evidência mínima:}

\begin{itemize}
  \item Alerta ZAP: Path Traversal
  \item OWASP: A01:2021 - Broken Access Control
  \item URL:
    \texttt{http://testasp.vulnweb.com/Templatize.asp?item=..\%2F..\%2F..\%2F..\%2F..\%2F..\%2F..\%2F..\%2F..\%2F..\%2F..\%2F..\%2F..\%2F..\%2F..\%2F..\%2FWindows\%2Fsystem.ini}
\end{itemize}

\section*{External Redirect}

\textbf{Exploração resumida:}

A aplicação permite redirecionamentos externos sem validação adequada da URL de
destino. Um atacante pode criar links maliciosos utilizando o domínio legítimo
da aplicação para redirecionar usuários para sites fraudulentos.

\textbf{Impacto provável:}

Phishing, roubo de credenciais, distribuição de malware e perda de confiança na
aplicação.

\textbf{Evidência mínima:}

\begin{itemize}
  \item Alerta ZAP: External Redirect
  \item OWASP: A01:2021 -- Broken Access Control
  \item URL:
    \texttt{http://testasp.vulnweb.com/Login.asp?RetURL=https\%3A\%2F\%2F8195536245220199150.owasp.org}
\end{itemize}

\newpage

\question[Atividade 3]{Associação das Metodologias de Pentest}

A atividade consistiu em associar as vulnerabilidades encontradas no
\textit{OWASP ZAP} com as categorias do \textit{OWASP Top 10 2021} e com as
táticas e técnicas do framework \textit{MITRE ATT\&CK}.

Inicialmente, foram analisadas as vulnerabilidades identificadas durante o
\textit{Automated Scan} realizado no site \texttt{http://testasp.vulnweb.com/}.

As vulnerabilidades selecionadas foram:

\begin{itemize}
  \item Cross Site Scripting (Reflected)
  \item SQL Injection
  \item Path Traversal
  \item External Redirect
\end{itemize}

Cada vulnerabilidade foi relacionada à categoria correspondente do
\textit{OWASP Top 10 2021}, considerando o tipo de risco apresentado pela
aplicação.

Em seguida, foi realizada a associação com o framework \textit{MITRE ATT\&CK},
identificando as táticas e técnicas mais prováveis utilizadas por atacantes
para explorar cada falha.

O mapeamento realizado foi:

\begin{itemize}
  \item \textbf{SQL Injection}
    \begin{itemize}
      \item OWASP: A03:2021 -- Injection
      \item Justificativa OWASP: a aplicação aceita entradas do usuário sem
        validação adequada, permitindo inserção de comandos SQL capazes de
        manipular consultas ao banco de dados e comprometer informações
        sensíveis.

      \item MITRE Tática: Initial Access
      \item MITRE Técnica: T1190 -- Exploit Public-Facing Application

      \item Justificativa MITRE: a vulnerabilidade pode ser explorada
        diretamente pela interface web exposta da aplicação, permitindo que um
        atacante obtenha acesso inicial por meio de requisições maliciosas
        identificadas pelo OWASP ZAP.
    \end{itemize}

  \item \textbf{Cross Site Scripting (Reflected)}
    \begin{itemize}
      \item OWASP: A03:2021 -- Injection
      \item Justificativa OWASP: entradas refletidas sem tratamento adequado
        permitem a execução de código JavaScript no navegador da vítima,
        podendo resultar em roubo de sessão ou manipulação do conteúdo exibido.

      \item MITRE Tática: Initial Access
      \item MITRE Técnica: T1189 -- Drive-by Compromise

      \item Justificativa MITRE: o ataque ocorre quando a vítima acessa uma URL
        manipulada contendo o script malicioso identificado nos alertas do ZAP,
        caracterizando execução automática de conteúdo malicioso no navegador.
    \end{itemize}

  \item \textbf{Path Traversal}
    \begin{itemize}
      \item OWASP: A01:2021 -- Broken Access Control
      \item Justificativa OWASP: a aplicação permite acesso inadequado a
        diretórios e arquivos internos, possibilitando que usuários acessem
        recursos restritos fora do diretório previsto pela aplicação.

      \item MITRE Tática: Initial Access
      \item MITRE Técnica: T1190 -- Exploit Public-Facing Application

      \item Justificativa MITRE: a exploração ocorre diretamente sobre a
        aplicação web pública, permitindo acesso indevido a arquivos do
        servidor através de parâmetros manipulados detectados pelo OWASP ZAP.
    \end{itemize}

  \item \textbf{External Redirect}
    \begin{itemize}
      \item OWASP: A01:2021 -- Broken Access Control
      \item Justificativa OWASP: a aplicação redireciona usuários para URLs
        externas sem validação adequada, possibilitando campanhas de phishing e
        envio da vítima para páginas maliciosas.

      \item MITRE Tática: Initial Access
      \item MITRE Técnica: T1189 -- Drive-by Compromise

      \item Justificativa MITRE: a vulnerabilidade pode induzir usuários a
        acessarem páginas controladas por atacantes, facilitando
        comprometimento inicial através de navegação maliciosa identificada
        durante a análise no ZAP.
    \end{itemize}
\end{itemize}

\newpage

Em forma de tabela temos:

\begin{table}[H]
  \centering
  \caption{Mapeamento entre vulnerabilidades, OWASP 2021 e MITRE ATT\&CK}
  \label{tab:owasp_mitre}
  \renewcommand{\arraystretch}{1.4}
  \begin{tabular}{|p{2.6cm}|p{1.9cm}|p{2.1cm}|p{2.5cm}|p{4.25cm}|}
    \hline
    \textbf{Vulnerabilidade (nome no ZAP)} &
    \textbf{Categoria OWASP 2021} &
    \textbf{MITRE ATT\&CK -- Tática} &
    \textbf{MITRE ATT\&CK -- Técnica (ID + nome)} &
    \textbf{Justificativa (1--2 linhas, baseada na evidência)} \\
    \hline

    SQL Injection &
    A03:2021 -- Injection &
    Initial Access &
    T1190 -- Exploit Public-Facing Application &
    A aplicação aceita entradas sem validação adequada, permitindo manipulação
    de consultas SQL através de parâmetros identificados nos alertas do OWASP
    ZAP. \\
    \hline

    Cross Site Scripting (Reflected) &
    A03:2021 -- Injection &
    Initial Access &
    T1189 -- Drive-by Compromise &
    Entradas refletidas sem sanitização permitem execução de JavaScript no
    navegador da vítima ao acessar URLs manipuladas identificadas pelo ZAP. \\
    \hline

    Path Traversal &
    A01:2021 -- Broken Access Control &
    Initial Access &
    T1190 -- Exploit Public-Facing Application &
    A aplicação permite acesso indevido a diretórios e arquivos internos por
    meio de parâmetros manipulados detectados durante a análise no OWASP ZAP.
    \\
    \hline

    External Redirect &
    A01:2021 -- Broken Access Control &
    Initial Access &
    T1189 -- Drive-by Compromise &
    A vulnerabilidade possibilita redirecionamento para páginas externas
    maliciosas, favorecendo ataques de phishing e comprometimento inicial do
    usuário. \\
    \hline

  \end{tabular}
\end{table}

Durante a análise, foi observado que a maior parte das vulnerabilidades estava
relacionada à fase de \textit{Initial Access}, indicando que as falhas
encontradas poderiam ser exploradas logo no primeiro contato do atacante com a
aplicação web.

Também foi possível identificar predominância de falhas relacionadas a injeção
e controle inadequado de acesso, demonstrando ausência de validação adequada
das entradas fornecidas pelo usuário.

As evidências utilizadas foram obtidas diretamente na seção \textit{Alertas} do
OWASP ZAP, incluindo nome do alerta, URL afetada e detalhes técnicos
apresentados pela ferramenta.

\newpage

\question[Atividade 4]{Análise de Resultados e Ações de Mitigação}

\textbf{Cross Site Scripting (Reflected)}

\begin{figure}[H]
    \centering
    \includegraphics[width=0.96\textwidth]{./fig/fig3.png}
    \caption{Detecção da vulnerabilidade de Cross Site Scripting refletido no OWASP ZAP.}
    \label{fig:xss}
\end{figure}

\begin{table}[H]
\centering
\begin{tabular}{|l|p{9cm}|}
\hline
\textbf{Risco} & Alto \\ \hline
\textbf{Confiança} & Média \\ \hline
\textbf{Parâmetro} & \texttt{tfText} \\ \hline
\textbf{Payload} & \texttt{</div><scrIpt>alert(1);</scRipt><div>} \\ \hline
\textbf{Evidência} & \texttt{</div><scrIpt>alert(1);</scRipt><div>} \\ \hline
\textbf{CWE} & 79 \\ \hline
\textbf{WASC} & 8 \\ \hline
\end{tabular}
\caption{Detalhes da vulnerabilidade Cross Site Scripting Refletido.}
\label{tab:xss}
\end{table}

\newpage
\textbf{External Redirect}

\begin{figure}[H]
    \centering
    \includegraphics[width=0.96\textwidth]{./fig/fig4.png}
    \caption{Redirecionamento externo identificado durante a análise.}
    \label{fig:redirect}
\end{figure}

\begin{table}[H]
\centering
\begin{tabular}{|l|p{9cm}|}
\hline
\textbf{Risco} & Alto \\ \hline
\textbf{Confiança} & Média \\ \hline
\textbf{Parâmetro} & \texttt{RetURL} \\ \hline
\textbf{Payload} & \texttt{https://8195536245220199150.owasp.org} \\ \hline
\textbf{Evidência} & \texttt{https://8195536245220199150.owasp.org} \\ \hline
\textbf{CWE} & 601 \\ \hline
\textbf{WASC} & 38 \\ \hline
\end{tabular}
\caption{Detalhes da vulnerabilidade External Redirect.}
\label{tab:redirect}
\end{table}

\newpage
\textbf{Path Traversal}

\begin{figure}[H]
    \centering
    \includegraphics[width=0.96\textwidth]{./fig/fig5.png}
    \caption{Exploração de diretórios através da vulnerabilidade Path Traversal.}
    \label{fig:pathtraversal}
\end{figure}

\begin{table}[H]
\centering
\begin{tabular}{|l|p{9cm}|}
\hline
\textbf{Risco} & Alto \\ \hline
\textbf{Confiança} & Média \\ \hline
\textbf{Parâmetro} & \texttt{item} \\ \hline
\textbf{Payload} & \texttt{../../../../../../../../Windows/system.ini} \\ \hline
\textbf{Evidência} & \texttt{[drivers]} \\ \hline
\textbf{CWE} & 22 \\ \hline
\textbf{WASC} & 33 \\ \hline
\end{tabular}
\caption{Detalhes da vulnerabilidade Path Traversal.}
\label{tab:pathtraversal}
\end{table}

\newpage
\textbf{SQL Injection}

\begin{figure}[H]
    \centering
    \includegraphics[width=0.96\textwidth]{./fig/fig6.png}
    \caption{Possível exploração de SQL Injection detectada pelo OWASP ZAP.}
    \label{fig:sqli}
\end{figure}

\begin{table}[H]
\centering
\begin{tabular}{|l|p{9cm}|}
\hline
\textbf{Risco} & Alto \\ \hline
\textbf{Confiança} & Baixa \\ \hline
\textbf{Parâmetro} & \texttt{tfSearch} \\ \hline
\textbf{Payload} & \texttt{ZAP'} \\ \hline
\textbf{Evidência} & \texttt{HTTP/1.1 500 Internal Server Error} \\ \hline
\textbf{CWE} & 89 \\ \hline
\textbf{WASC} & 19 \\ \hline
\end{tabular}
\caption{Detalhes da vulnerabilidade SQL Injection.}
\label{tab:sqli}
\end{table}

\newpage
\textbf{Absence of Anti-CSRF Tokens}

\begin{figure}[H]
    \centering
    \includegraphics[width=0.96\textwidth]{./fig/fig7.png}
    \caption{Formulário sem proteção Anti-CSRF identificado na aplicação.}
    \label{fig:csrf}
\end{figure}

\begin{table}[H]
\centering
\begin{tabular}{|l|p{9cm}|}
\hline
\textbf{Risco} & Médio \\ \hline
\textbf{Confiança} & Baixa \\ \hline
\textbf{Evidência} & \texttt{<form action="" method="post"...>} \\ \hline
\textbf{CWE} & 352 \\ \hline
\textbf{WASC} & 9 \\ \hline
\end{tabular}
\caption{Ausência de tokens Anti-CSRF.}
\label{tab:csrf}
\end{table}

\newpage
\textbf{Content Security Policy (CSP) Header Not Set}

\begin{figure}[H]
    \centering
    \includegraphics[width=0.96\textwidth]{./fig/fig8.png}
    \caption{Ausência do cabeçalho CSP identificada pelo OWASP ZAP.}
    \label{fig:csp}
\end{figure}

\begin{table}[H]
\centering
\begin{tabular}{|l|p{9cm}|}
\hline
\textbf{Risco} & Médio \\ \hline
\textbf{Confiança} & Alta \\ \hline
\textbf{CWE} & 693 \\ \hline
\textbf{WASC} & 15 \\ \hline
\end{tabular}
\caption{Ausência do cabeçalho Content Security Policy.}
\label{tab:csp}
\end{table}

\newpage
\textbf{Missing Anti-clickjacking Header}

\begin{figure}[H]
    \centering
    \includegraphics[width=0.96\textwidth]{./fig/fig9.png}
    \caption{Cabeçalhos de proteção contra Clickjacking ausentes.}
    \label{fig:clickjacking}
\end{figure}

\begin{table}[H]
\centering
\begin{tabular}{|l|p{9cm}|}
\hline
\textbf{Risco} & Médio \\ \hline
\textbf{Confiança} & Média \\ \hline
\textbf{CWE} & 693 \\ \hline
\textbf{WASC} & 15 \\ \hline
\end{tabular}
\caption{Ausência de proteção contra Clickjacking.}
\label{tab:clickjacking}
\end{table}

\begin{table}[H]
\centering
\small
\begin{tabular}{|p{5cm}|c|c|c|}
\hline
\textbf{Vulnerabilidade} & \textbf{Risco} & \textbf{CWE} & \textbf{WASC} \\ \hline
Application Error Disclosure & Baixo & 550 & 13 \\ \hline
Cookie No HttpOnly Flag & Baixo & 1004 & 13 \\ \hline
Cookie without SameSite Attribute & Baixo & 1275 & 13 \\ \hline
Information Disclosure - Debug Error Messages & Baixo & 1295 & 13 \\ \hline
Server Leaks Information via X-Powered-By & Baixo & 497 & 13 \\ \hline
Server Leaks Version Information via Server Header & Baixo & 497 & 13 \\ \hline
X-Content-Type-Options Header Missing & Baixo & 693 & 15 \\ \hline
Session Management Response Identified & Informativo & -- & -- \\ \hline
User Controllable HTML Element Attribute & Informativo & 20 & 20 \\ \hline
\end{tabular}
\caption{Resumo das vulnerabilidades de baixo risco e informacionais.}
\label{tab:lowrisk}
\end{table}
